\documentclass[aspectratio=169]{beamer}

\usepackage[utf8]{inputenc}
\usepackage[T1]{fontenc}
\usepackage[portuguese]{babel}
\usepackage{helvet}
\usepackage{enumitem}
\usepackage{xcolor}
\usepackage{hyperref}
\usepackage{array}

\renewcommand{\familydefault}{\sfdefault}

\definecolor{iselblue}{RGB}{45,62,230}
\definecolor{textgray}{RGB}{55,55,55}

\setbeamercolor{normal text}{fg=textgray,bg=white}
\setbeamercolor{structure}{fg=iselblue}
\setbeamercolor{frametitle}{fg=iselblue,bg=white}
\setbeamercolor{title}{fg=iselblue}
\setbeamercolor{footline}{fg=iselblue,bg=white}
\setbeamercolor{item}{fg=black}

\setbeamertemplate{navigation symbols}{}
\setbeamertemplate{items}[circle]
\setbeamertemplate{footline}{
  \leavevmode%
  \hbox{%
    \begin{beamercolorbox}[wd=.15\paperwidth,ht=2.6ex,dp=1.2ex,leftskip=1em]{footline}%
      LEIC
    \end{beamercolorbox}%
    \begin{beamercolorbox}[wd=.70\paperwidth,ht=2.6ex,dp=1.2ex,center]{footline}%
      ShopISEL
    \end{beamercolorbox}%
    \begin{beamercolorbox}[wd=.15\paperwidth,ht=2.6ex,dp=1.2ex,rightskip=1em,right]{footline}%
      PS 2024/2025
    \end{beamercolorbox}%
  }%
  \vskip0pt%
}
\setbeamertemplate{headline}{
  \vspace{0.45cm}
  \hspace{0.08\paperwidth}\color{iselblue}\rule{0.84\paperwidth}{0.5pt}
}

\title{Guia para o Arguente}
\subtitle{ShopISEL}
\author{Martim Monteiro \and Alexandre Tavares \and Pedro Gomes}
\date{Julho de 2026}

\begin{document}

\begin{frame}
  \frametitle{Descri\c{c}\~ao da organiza\c{c}\~ao}
  \begin{itemize}[leftmargin=1.6em]
    \item \textbf{T\'itulo do projeto:} ShopISEL -- Aplica\c{c}\~ao de Lista de Compras com Pre\c{c}os e Alertas
    \item \textbf{Grupo:} Martim Monteiro (51619), Alexandre Tavares (51271), Pedro Gomes (51423)
    \item \textbf{Orientador:} Doutor Fernando Miguel Gamboa de Carvalho
    \item \textbf{Data:} Julho de 2026
    \item \textbf{Objetivo desta sec\c{c}\~ao:} dar ao arguente um ponto de entrada r\'apido para consultar, executar e validar a solu\c{c}\~ao.
  \end{itemize}
\end{frame}

\begin{frame}
  \frametitle{Organiza\c{c}\~ao do dossier do projeto}
  \small
  \begin{itemize}[leftmargin=1.6em]
    \item \textbf{Frontend web:} \texttt{frontend/} -- aplica\c{c}\~ao React com Vite.
    \item \textbf{Frontend mobile:} \texttt{frontend-mobile/} -- aplica\c{c}\~ao Expo React Native.
    \item \textbf{Servi\c{c}os backend:} \texttt{account-service/}, \texttt{list-service/}, \texttt{products-service/}, \texttt{prices-service/}, \texttt{stores-service/}, \texttt{matchqueue-service/}.
    \item \textbf{Recolha de dados:} \texttt{scraper-continente/}, \texttt{scraper-pingodoce/}, \texttt{all-stores-scraper/}, \texttt{scraper-orchestrator/}.
    \item \textbf{Autentica\c{c}\~ao:} \texttt{keycloack/} com configura\c{c}\~ao de deploy do Keycloak.
    \item \textbf{Gateway/deploy:} \texttt{deploy/} com configura\c{c}\~oes Fly.io e Render.
    \item \textbf{Relat\'orio:} \texttt{report/report/}.
    \item \textbf{Acesso ao c\'odigo:} disponibilizar o URL do reposit\'orio do grupo e permiss\~oes de leitura ao arguente antes da prova.
  \end{itemize}
\end{frame}

\begin{frame}[fragile]
  \frametitle{Como instalar e executar}
  \small
  \begin{itemize}[leftmargin=1.6em]
    \item \textbf{Web frontend}
\begin{verbatim}
cd frontend
npm install
npm run dev
\end{verbatim}
    \item \textbf{Mobile frontend}
\begin{verbatim}
cd frontend-mobile
npm install
npx expo start --go
\end{verbatim}
    \item \textbf{Servi\c{c}os .NET}
\begin{verbatim}
cd <nome-do-servico>
dotnet restore src
dotnet test src
dotnet run --project src/<ServiceName>
\end{verbatim}
    \item \textbf{Configura\c{c}\~ao:} preencher vari\'aveis de ambiente e \texttt{ConnectionStrings} dos servi\c{c}os; no frontend web usar \texttt{frontend/.env}; no Keycloak definir os segredos descritos em \texttt{keycloack/README.md}.
  \end{itemize}
\end{frame}

\begin{frame}
  \frametitle{Acessos e guia de valida\c{c}\~ao}
  \small
  \begin{itemize}[leftmargin=1.6em]
    \item \textbf{Credenciais de acesso:} inserir aqui a conta de demonstra\c{c}\~ao do Keycloak e, se aplic\'avel, conta de administrador do realm.
    \item \textbf{Exemplo:} utilizador \texttt{[preencher]} e palavra-passe \texttt{[preencher]}.
    \item \textbf{Servi\c{c}os externos necess\'arios:} base de dados PostgreSQL, Keycloak e Firebase Cloud Messaging.
    \item \textbf{Verifica\c{c}\~oes sugeridas ao arguente:}
    \begin{itemize}[leftmargin=1.6em]
      \item autenticar no sistema;
      \item criar e editar uma lista de compras;
      \item consultar produtos e pre\c{c}os;
      \item confirmar integra\c{c}\~ao entre frontend, APIs e autentica\c{c}\~ao;
      \item rever a organiza\c{c}\~ao dos microservi\c{c}os e pipelines de deploy.
    \end{itemize}
    \item \textbf{Nota:} esta slide deve ser finalizada com o URL real do reposit\'orio, contas de demonstra\c{c}\~ao e eventuais passos de arranque espec\'ificos do ambiente da prova.
  \end{itemize}
\end{frame}

\end{document}
